\section{Comparative Computation}\label{sec:r45-study}
The study separates exact conditional verification, joint optimization quality, independent model validation, and scaling. Every input, rational policy, seed, budget frontier, numerical solver result, and repeated timing is retained. The experiments are synthetic and make no population or field-performance claim. Tables~\ref{tab:r45-comparison}--\ref{tab:r45-continuous} summarize executed files generated by the accompanying scripts, not manually entered performance claims.

\subsection{Conditional correctness and the global target net}
Six declared seeds produce 144 small instances with distinct and repeated caps, heterogeneous costs, several risk tolerances, and nonnegative charges. For each, an independent routine enumerates catalog subsets and all singleton and two-point supports at the supplied targets. It does not import the new recurrence, the old price decomposition, or the canonical reconstruction. All 288 tested budget values agree exactly with the conditional dynamic program. The suite also checks 144 target repairs, fifteen comparison-book rate cases, four endpoint cases, ten semantic corruption rejections, and five invalid-input rejections. These are regression and theorem checks, not optimization-performance evidence.

The target-net scheme is tested separately on twenty two- and three-branch settings, including nonuniform charges. Each completed run is compared with an exact global catalog optimum obtained by enumerating books and solving their concave inner allocations. The actual loss is no larger than the requested additive tolerance in every run; the largest observed loss is $0.0498046875$. Together these runs evaluate 385 feasible profiles, with at most 94 in a run. This modest experiment verifies the net construction and its coverage certificate; it does not disguise the exponential dependence on the number of branches. An interrupted run explicitly returns no global accuracy certificate.

\subsection{Original and enlarged budget comparisons}
The comparative sample has six seeds, three cap families (spread, clustered, and repeated), and both free and charged interfaces, for 36 instances. It varies weights, shortfall curvatures, risk tolerance, promise saturation, and arbitrary nonnegative charges. Each instance has three or four branches, seven or nine catalog levels, and original budget two or three. Exact catalog optima are computed at every budget through twice the original budget. This gives the original optimum, the optimum at the actual union size, and the full enlarged-budget reference rather than treating an original-budget price bound as an enlarged-budget optimality gap.

The core candidate method uses three feasible target profiles: capped ideal targets, the two-price mixed targets, and the original-budget incumbent's targets. For each profile it computes the entire conditional frontier, then uses at most twelve monotone allocation--book improvement steps per budget. This finite target portfolio is a heuristic for joint design, not a substitute for the completed target net. Baselines are the old two-price recovery, greedy level addition with exact inner reallocation, greedy deletion of union levels with reallocation, and direct exact catalog optimization. The operational recommendation retains the better of the core method and the greedy policy at each budget. All component policies remain separately reported.

The core method attains the original-budget exact optimum on all eighteen free instances and seventeen of eighteen charged instances. It improves on greedy addition in ten cases; greedy addition is better in one. The charged exception, seed 205019 in the spread family, is important: the core target portfolio is $101/3600$ below the optimum $1373/3600$, while the greedy policy attains that optimum. The safeguard therefore attains all 36 original-budget optima in this sample. This observation is not a universal guarantee for the safeguard, nor is the exception removed from the data. It demonstrates why conditional global optimality must not be presented as global target optimality.

Two instances improve strictly on the old original-budget incumbent. Across all instances, endpoint books are nested in 23 cases; the largest observed union has four levels. The data retain actual union size, overlap, nonnegative expansion $\max\{0,|c^U|-m\}$, percentage expansion, each budget's joint reference, and union loss at its own size. These outcomes do not establish that a factor-two expansion is necessary. The new conditional frontier instead answers how much each permitted budget can recover for the supplied targets, and the exact small references show the remaining target-selection loss.

\subsection{A genuinely different optimization formulation}
For twelve selected input instances, we solve the original uneliminated model at two budgets, giving 24 comparisons and 48 mixed-integer envelope solves. The model retains terminal probabilities, binary level activations, binary branch supports, a pre-draw intermediate tier, the exact root equality, and the implication that every positive-probability realization satisfies $y_j+a_i\leq\tau_j$. It contains neither target-crossing recurrence nor the canonical response formula. Section~\ref{ec:r45-mip} gives the complete formulation.

Quadratic costs are bounded by tangent and secant piecewise-linear envelopes. The former gives a maximization upper problem; the latter gives a conservative feasible lower problem. Each uses 32 segments per branch, with an analytic envelope-error allowance. SciPy's mixed-integer interface and HiGHS solve these models in floating-point arithmetic \citep{Virtanen2020,HuangfuHall2018}. In all 24 comparisons the exact catalog value lies between the numerical brackets within the declared $10^{-7}$ comparison tolerance, both solves terminate with reported optimal status, and direct model residuals are below $10^{-6}$. The tangent-to-secant width is at most twice the envelope allowance plus $2\times10^{-6}$. These are independently formulated \emph{numerical} checks, not exact rational proofs of global optimality. The exact comparisons and rational policy checks supply different evidence and are labeled accordingly.

\subsection{Scaling, local certificates, and continuous references}
Twelve further cases use two seeds, branch counts 32, 128, and 256, catalog sizes 17, 33, and 65, and budgets four or eight. Unlike the earlier deterministic family, caps, weights, curvatures, charges, and branch-specific realization headroom vary randomly. Each computes the ideal-target conditional frontier three times. Table~\ref{tab:r45-scaling} reports the median of these timings and the same-budget Jensen gap. These runs establish conditional scalability and feasible value certificates, not global optimization of a 256-dimensional target space. Timing replicates measure implementation variability on one machine and are not independent draws of an applied population. Stored dynamic-programming bit lengths exclude unrecorded arithmetic temporaries. The archived process peak memory is cumulative operating-system resident memory, not a sub-megabyte allocation-tracing statistic.

The continuous comparison uses four explicitly solvable regimes. In an interior two-branch case, caps $(3/10,9/10)$, weights $(2/5,3/5)$, zero headroom, and promise $66/125$ give ideal targets $(3/10,17/25)$. A four-candidate adaptive pool containing these targets attains the continuous two-symbol optimum exactly. The uniform 65-point catalog still has a loss of $0.0029909375$. For a saturated off-cap-top example, midpoint refinement reaches the exact continuous optimum with six candidates, whereas the nine-point uniform pool also reaches it. Two additional seven-branch saturated cases use the retained exact continuous optimizer as reference; target-adaptive pools with nine candidates attain their optima, while uniform pools require 33 points among the tested resolutions. All these comparisons keep the installed-symbol budget fixed. They illustrate the advantage of aligning candidate points with the problem, not a universal rate for adaptive refinement. Their exact-reference errors are distinct from the Jensen certificate, which can include intrinsic memory loss.

Finally, ten deliberately suboptimal support-policy pairs verify the inexact-price identity with positive oracle defects, rather than merely retest exact optimizers. The target-repair tests reconstruct rational lotteries and check realized totals directly. Neither exercise validates unknown estimated model parameters. For uncertain caps or realization limits, a deployment must first supply conservative feasible limits; arbitrarily perturbing an eligibility threshold can change the feasible support discretely.
